\section{Computer-Based Interlocking (CBI)}

Computer-Based Interlocking (CBI) represents a fundamental shift from hardware-defined safety logic, such as relays or mechanical linkages, toward software-defined, programmable control. In a CBI, interlocking rules are expressed in software, often using high-level programming languages or specialized safety-certified platforms. This approach offers several advantages compared to traditional systems. Configurability becomes a key strength: rules can be updated, extended, or debugged without requiring physical rewiring or hardware replacement, significantly reducing cost and downtime. Scalability is also improved, as the same software framework can be adapted to different station layouts, operational policies, or national signaling practices. Another major benefit is the possibility of simulation-driven development—scenarios can be modeled and tested in a digital environment before deployment, allowing potential errors to be identified early. Finally, CBIs are inherently integration-ready, capable of interfacing with SCADA systems, traffic control centers, and modern IoT-based monitoring infrastructure, thereby enabling railways to operate as part of a larger, digitized transport ecosystem.

However, these advantages are balanced by significant challenges. The most critical is safety assurance, since shifting safety logic into software introduces risks of programming errors, unexpected interactions, or runtime faults. Unlike hardware logic, which is deterministic by design, software must be validated through rigorous verification pipelines, including formal proofs, redundancy strategies, and extensive field testing. Hardware dependency also poses risks, as failures in execution environments—such as operating system crashes, memory corruption, or processor faults—can compromise safety if not properly mitigated. Furthermore, CBIs must comply with demanding international standards, including EN 50128 for railway software and IEC 61508 for functional safety, which mandate strict processes for development, testing, and certification. Achieving compliance requires specialized expertise, documentation, and continuous auditing, often increasing both project cost and development time.

Despite these challenges, the global trend toward CBIs is accelerating, especially in regions seeking to modernize aging infrastructure without complete hardware overhaul. By combining flexibility, scalability, and digital integration with safety-critical assurance frameworks, CBIs are becoming the cornerstone of next-generation railway signaling.